\section{Keep lease rights distinct within one company book}
\label{sec:lease-shared-book}

Preserving a proposal makes its calculations available for inspection. To
reproduce them, the company model must also preserve the rights that generated
the entries. A total lease liability and a total right-of-use asset do not tell
us how much of either balance disappears when one right ends. Nor does a
balanced calculation for one lease establish that the company can pay it after
another lease has used the same cash.

Consider the retained synthetic example with two rights in one lease. Right A
has a right-of-use asset of 90 and a liability of 100; right B has an asset of
30 and a liability of 50. Amounts are exact USD quantities in this example.
The supplied findings support tracing the amounts to their individual rights.
Terminating B therefore removes its asset of 30 and liability of 50, giving a
gain of 20. Allocating the total asset of 120 in proportion to B's share of the
total liability would instead remove
\[
  120\frac{50}{100+50}=40,
  \qquad 50-40=10.
\]
That gain of 10 is wrong for the stated traced-right inputs. The two rights
have different asset-to-liability ratios. Their totals conceal that difference,
so the simulator retains each right's amounts through termination and
remeasurement.

This example uses the inspected \path{CorporateResiduals/1} methods and their
historical synthetic source scope. The owner remains
\path{Synthetic-heterogeneous-rights}, with the Legal book, base branch, USD
currency and \path{historical_EU_IFRS_conditional} basis. The permitted execution
window begins on 1 January 2024 and ends immediately before 1 January 2026.
Complete opening balances, lease membership and method findings are supplied
premises. They are not inferred from an actual issuer's reports or evidence of
current accounting compliance.

\subsection{Follow the rights through the complete lifecycle}

Both opening rights have a supplied annual rate of 10 percent. On 1 January
2024, the example terminates B and replaces A's payment schedule with one
payment of 121 at the annual period end. The original method discounts that
payment for one year. For A, the liability rises from 100 to 110, and the
remeasurement adds the same 10 to its asset:
\[
  L_A^{(1)}=\frac{121}{1.1}=110,
  \qquad R_A^{(1)}=90+(110-100)=100.
\]
Here $L_A$ is A's positive liability amount and $R_A$ its right-of-use asset;
the superscript identifies the completed step. B is now zero. Combining B's
termination with A's remeasurement reduces the total asset by 20 and the
positive liability by 40. With debit-positive signed accounts, the journal is
therefore $\Delta\mathrm{rou}=-20$,
$\Delta\mathrm{lease\_liability}=+40$ and
$\Delta\mathrm{other\_gains}=-20$. The entries sum to zero; the credit to
gains records the 20 profit.

A second supplied modification occurs on the same date. It removes a further
10 of A's asset and 11 of its liability, producing a gain of 1. The remaining
amounts before remeasurement are asset 90 and liability 99. Replacing the
remaining annual payment with 99 gives:
\[
  L_A^{(2)}=\frac{99}{1.1}=90,
  \qquad R_A^{(2)}=(100-10)+(90-(110-11))=81.
\]
The remeasurement reduces the remaining liability by 9 and makes the same
reduction to the asset. Across this second step, the total asset falls by 19,
the positive liability falls by 20, and the additional gain is 1. Tracing the
removed amounts and then remeasuring the remainder explains all three changes.

On 1 January 2025, the supplied findings confirm one elapsed year with no
intervening accrual or modification. The payment of 99 comprises interest of
$90\times0.1=9$ and principal of $99-9=90$. Cash falls by 99, finance expense
increases by 9, and the signed liability increases by 90 to zero. The retained
cash-flow policy assigns 9 to operating outflow and 90 to financing outflow.
Finally, a supplied systematic depreciation amount of 81 removes the remaining
asset and records expense of 81. This depreciation amount is an input supported
by the example's finding; the wrapper does not infer a useful life.

With opening cash 200, closing cash is $200-99=101$. Both lease balances are
zero, and cumulative pretax profit from these four events is
\[
  \Pi=20+1-9-81=-69.
\]
The gains, interest and depreciation explain the profit independently of the
cash payment. This is a complete lifecycle for the supplied finite example;
it does not supply an annual tax close or a complete financial report.

\subsection{Apply each lease to the same cash balance}

Now instantiate a second synthetic lease with the same supplied right-level
amounts and its own identity and premise. A company opening containing both
leases must reconcile to assets of 240 and positive liabilities of 300. Cash
and equity belong to that one company opening. Adding an isolated lease
calculator's opening cash or equity would capitalize the company again each
time a method runs.

For example, opening company cash of 400 followed by one payment of 99 gives
301. Adding the component fixture's opening cash of 200 would give 501, an
incorrect result. The shared-book proposal instead applies only the original
journal deltas. It updates the selected lease's rights, preserves the other
lease's memory, and leaves unrelated company accounts unchanged. Right A in
the first lease and right A in the second lease are distinct: their identities
are the pairs of lease identifier and local right identifier.

Funding must be checked in event order. Under the explicitly supplied
\path{no_overdraft} policy, company cash of 300 funds both payments of 99 and
leaves 102. Cash of 150 funds the first payment and leaves 51. The second
payment then fails: its liability remains 90, while the first lease's liability
remains settled. The successful earlier events survive; the failed event adds
no journal, cash movement or after-state. This is a funding failure under the
company premise. The original component calculator's ability to calculate a
balanced entry ending in negative cash does not provide that funding policy.

\subsection{Bind proposals to complete state and original methods}

The \path{SharedLeaseBook/1} record contains context, original source-scope
hash, execution period, current date, the full original account registry, all
lease components, collection-completeness and instantiation premises, funding
policy, prior event identities, gross cash movements and the conditional-state
flag. Every right retains its asset, liability, rate and last measurement date.
Canonical rational strings preserve exact amounts. Before and after each event,
the signed balances sum to zero, the ROU account equals all named rights'
assets, and the negative lease-liability account equals all their liabilities.
Unknown membership cannot establish those equalities merely because a supplied
subset balances.

\path{LeaseBookEvent/1} binds a new event identity, lease identity, date, source
scope and one of the original \path{lease.reduce}, \path{lease.pay} or
\path{lease.depreciate} methods. Parameters must satisfy that method's closed
schema and required findings. Dates are nondecreasing within the execution
window; a selected right's measurement date cannot lie after the event.
Payment retains the original exact annual-anniversary requirement. A missing
required date does not create a dated prefix, and no leap-day convention is
silently supplied.

The permitted account changes are explicit. Reduction affects ROU, lease
liability and other gains. Payment affects cash, finance expense and lease
liability. Depreciation affects ROU and depreciation expense. The wrapper checks
these effects, reconciles the complete rights collection and enforces the
company's cash policy. Unknown supported inputs remain unresolved; malformed
records, false findings and unsupported schedules retain their distinct original
refusal categories. No invented loan or balancing entry repairs a refusal.

A \path{LeaseBookProposal/1} binds the complete before-state, event and original
method hashes to the hypothetical after-state, journal and gross cash effects.
Applying it recomputes the original method and compares the entire proposal.
A forged cash increase offset by equity can preserve balance while failing
this comparison. Context, source, period, unrelated accounts and nonselected
leases must remain unchanged. The result remains hypothetical:
\path{committed}, \path{source_admission}, \path{accounting_admission} and
\path{method_eligible} are false. Application here does not publish a database
transaction; an evidence-archive receipt preserves bytes without replacing
financial recomputation.

The original-scope audit and shared-book reference, including complete input
and output examples, are retained under \path{data/accounting-master/}:
\begin{quote}\small
\path{lease-module-composition-audit-20261024-r2}\\
\path{lease-shared-book-20261024}
\end{quote}
The checked composition preserves heterogeneous rights and one cash balance
for these methods. General lease valuation, FX, tax and report integration,
arbitrary module imports, and reconstruction of actual contract rights remain
separate design dependencies. The following joined-company profile addresses
another finite combination of legal books and methods under its own premises.
